% TOP_PROBLEM: {"id": 3185, "title": "Seagull problem", "category_id": 3, "category": {"id": 3, "name": "graph_theory", "display_name": "Graph Theory", "description": "Problems involving graphs, networks, and their properties.", "slug": "graph-theory", "order_index": 3, "created_at": "2026-07-31T15:26:25.671Z"}, "status": "open", "problem_number": "OPG-729", "set_id": 12}
% FIRST_POSTED: 2026-10-01
% ATTEMPT_STATUS: unresolved
% UnsolvedMath ID 3185; source catalogue code OPG-729
% !TeX program = lualatex
% GPT-6 Astra Ultra research attempt, 2026-10-01; independently checked partial work.
% Completion estimate: no numerical percentage assigned; see final status and literature audit.
\documentclass[11pt,a4paper]{article}
\usepackage[margin=25mm]{geometry}
\usepackage{fontspec}
\setmainfont{lmroman10-regular.otf}[BoldFont=lmroman10-bold.otf,
 ItalicFont=lmroman10-italic.otf,BoldItalicFont=lmroman10-bolditalic.otf]
\usepackage{amsmath,amssymb,mathtools}
\usepackage{unicode-math}
\setmathfont{latinmodern-math.otf}
\AtBeginDocument{\let\setminus\smallsetminus}
\usepackage{xurl,hyperref}
\hypersetup{colorlinks=true,urlcolor=blue}
\setcounter{secnumdepth}{0}
\setlength{\parindent}{0pt}
\setlength{\parskip}{5pt}
\setlength{\emergencystretch}{3em}
\begin{document}
\section*{Seagull problem}\label{3185}
{\small UnsolvedMath ID 3185 · OPG-729 · Graph Theory · OpenGarden}
\subsection{Definitions and mathematical statement}
Let $G=(V,E)$ be a finite simple graph with
$n=|V|\geq1$. An independent set is a subset of vertices
with no edges between any two of them. The hypothesis
that there is no independent set of size three is
equivalent to saying that among every three distinct
vertices some pair is adjacent, or that the independence
number $\alpha(G)$ is at most two.

For an integer $t\geq1$, a $K_t$ minor can be specified
by pairwise disjoint nonempty branch sets
$B_1,\ldots,B_t\subseteq V$, such that each $G[B_i]$
is connected and there is an edge between every pair
of distinct branch sets. Contracting each branch set
and deleting surplus vertices and edges produces the
complete graph $K_t$.

The conjecture is
\[
 \alpha(G)\leq2\quad\Longrightarrow\quad
 G\text{ has a }K_{\lceil n/2\rceil}\text{ minor}.
\]
The ceiling gives the precise integer interpretation
of a complete minor on at least $n/2$ vertices.
No connectivity assumption is imposed on $G$, and
the branch sets need not exhaust its vertices.

The name refers to a seagull, an induced three-vertex
path: its two ends are nonadjacent and its middle
vertex is adjacent to both. Such connected triples
can serve as branch sets in attempts to build the
required minor, but the conjecture places no shape
or size restriction on its branch sets. It asks
for a complete minor, not just a dense minor or
a large clique subgraph.
\subsection{Short English statement}
Must every n-vertex graph with no three pairwise nonadjacent vertices have a complete minor on at least half its vertices?
\subsection{Research attempt}
\paragraph{Scope and process.}
GPT-6 Astra Ultra, 1 October 2026. The catalogue formulation is retained above. This attempt checks the original problem and primary literature, proves the restricted claims stated below, and identifies the exact limit of its conclusions. Reproved facts are not claimed as novel results.

\paragraph{Literature and attack.}
The original Garden entry [S2] explains the seagull contraction argument and cites stronger refinements. We do not treat its historical discussion as a current best-bound claim. Here we audit the branch-set accounting, extract a precise sufficient condition for reaching the catalogue threshold, and give an explicit family where the accounting loses vertices despite an easily visible target minor. A seagull means an induced three-vertex path, not an arbitrary connected triple.

\paragraph{Basic certificate, with all adjacency conditions.}
Choose a maximal collection of disjoint seagulls, say $S_1,\ldots,S_t$. Every vertex outside a seagull has a neighbour among its two nonadjacent endpoints, since otherwise those three vertices would be independent. Thus each $S_i$ is adjacent to every other chosen branch set. The residual graph has no induced three-vertex path, so its components are cliques: a shortest path between nonadjacent vertices in a connected component would supply such a path. There are at most two residual components, because one vertex in each of three components would be an independent triple. Write their orders as $a\ge b\ge0$. Contracting the $t$ seagulls and taking every vertex of the larger clique gives a complete minor of order $r=t+a$. These statements justify connectivity, disjointness and pairwise adjacency of the branch sets, not only their count.

\paragraph{Exact deficit and a restricted theorem.}
Since $n=3t+a+b$, the above certificate obeys
\[
 2r-n=a-b-t. \tag{1}
\]
It reaches $r\ge\lceil n/2\rceil$ exactly when $a-b\ge t$. Therefore the catalogue assertion is proved for any graph admitting a maximal seagull packing with residual clique imbalance at least the number of seagulls. This criterion is testable from the packing itself and does not require it to have maximum cardinality. A maximal packing means simply that another disjoint induced path cannot be added; that is precisely what ensures the residual structure used above.

For example, if no seagull exists, then $t=0$ and $G$ consists of at most two cliques, so the larger clique already has the required order. If one seagull has been removed and the residual clique orders differ by at least one, the same theorem gives the target. The integer rounding is exact: (1) has the same parity as $n$, and $2r\ge n$ is equivalent to $r\ge\lceil n/2\rceil$.

The unconditional count obtained from this certificate is
\[
 r=t+a\ge t+\frac{a+b}{2}=\frac{n-t}{2}\ge\frac n3,
\]
since $3t\le n$. This is only the elementary baseline, not an assertion of optimality among known results.

\paragraph{An explicit limitation of arbitrary maximal packings.}
For any $t\ge1$, take $t$ disjoint triples $\{x_i,y_i,z_i\}$, include $x_i y_i,y_i z_i$ but omit $x_i z_i$, and add every edge between different triples. Call the resulting graph $G_t$. Its only nonedges are the pairs $x_i z_i$, so $\alpha(G_t)=2$. The displayed triples form a maximal packing of $t$ seagulls exhausting all vertices. For this packing $a=b=0$, and the certificate above has only $t$ branch sets, whereas the target is $\lceil3t/2\rceil$.

The graph itself is no obstruction: choosing $x_i,y_i$ from every triple gives a clique on $2t$ vertices, which exceeds or equals the target. Thus even a packing of maximum possible size can be a poor guide to the desired minor. Optimizing the number of seagulls alone is the wrong objective for this method. The example isolates an actual failure of the proposed accounting without masquerading as a counterexample to the conjecture.

\paragraph{Verification and remaining gap.}
The script \texttt{scripts/check\_attempt\_batch02\_graphs\_2026\_10\_01.py} checks this family for $1\le t\le6$, including independence number two, the induced-path partition and the $2t$-clique. The symbolic construction proves it for all $t$. A general proof would need either a packing with $a-b\ge t$ or a way to reuse vertices from seagulls as additional compatible branch sets. Neither property follows from maximality, and no such universal improvement is established here.

\paragraph{Final status.}
The residual-imbalance theorem and the explicit failure of the packing-count heuristic are proved. The full $K_{\lceil n/2\rceil}$-minor target is \textbf{UNRESOLVED} by this attempt.

\subsection{Sources}
\begin{itemize}
\item[S1] ULAM.AI and the UnsolvedMath Contributors, supplied problems.json, record 3185 (OPG-729), OpenGarden collection. Catalogue identity and source transcription; CC BY 4.0.
\url{https://huggingface.co/datasets/ulamai/UnsolvedMath}.
\item[S2] Open Problem Garden, Seagull problem. Original problem and discussion; the mathematical formulation below is expanded from the supplied source record.
\url{http://www.openproblemgarden.org/op/seagull_problem}.

\end{itemize}
\end{document}
